% Options for packages loaded elsewhere
\PassOptionsToPackage{unicode}{hyperref}
\PassOptionsToPackage{hyphens}{url}
%
\documentclass[
]{ctexart}
\usepackage{amsmath,amssymb}
\usepackage{lmodern}
\usepackage{iftex}
\ifPDFTeX
  \usepackage[T1]{fontenc}
  \usepackage[utf8]{inputenc}
  \usepackage{textcomp} % provide euro and other symbols
\else % if luatex or xetex
  \usepackage{unicode-math}
  \defaultfontfeatures{Scale=MatchLowercase}
  \defaultfontfeatures[\rmfamily]{Ligatures=TeX,Scale=1}
\fi
% Use upquote if available, for straight quotes in verbatim environments
\IfFileExists{upquote.sty}{\usepackage{upquote}}{}
\IfFileExists{microtype.sty}{% use microtype if available
  \usepackage[]{microtype}
  \UseMicrotypeSet[protrusion]{basicmath} % disable protrusion for tt fonts
}{}
\makeatletter
\@ifundefined{KOMAClassName}{% if non-KOMA class
  \IfFileExists{parskip.sty}{%
    \usepackage{parskip}
  }{% else
    \setlength{\parindent}{0pt}
    \setlength{\parskip}{6pt plus 2pt minus 1pt}}
}{% if KOMA class
  \KOMAoptions{parskip=half}}
\makeatother
\usepackage{xcolor}
\IfFileExists{xurl.sty}{\usepackage{xurl}}{} % add URL line breaks if available
\IfFileExists{bookmark.sty}{\usepackage{bookmark}}{\usepackage{hyperref}}
\hypersetup{
  hidelinks,
  pdfcreator={LaTeX via pandoc}}
\urlstyle{same} % disable monospaced font for URLs
\usepackage[margin=2.5cm]{geometry}
\usepackage{longtable,booktabs,array}
\usepackage{calc} % for calculating minipage widths
% Correct order of tables after \paragraph or \subparagraph
\usepackage{etoolbox}
\makeatletter
\patchcmd\longtable{\par}{\if@noskipsec\mbox{}\fi\par}{}{}
\makeatother
% Allow footnotes in longtable head/foot
\IfFileExists{footnotehyper.sty}{\usepackage{footnotehyper}}{\usepackage{footnote}}
\makesavenoteenv{longtable}
\setlength{\emergencystretch}{3em} % prevent overfull lines
\providecommand{\tightlist}{%
  \setlength{\itemsep}{0pt}\setlength{\parskip}{0pt}}
\setcounter{secnumdepth}{-\maxdimen} % remove section numbering
\ifLuaTeX
  \usepackage{selnolig}  % disable illegal ligatures
\fi

\author{}
\date{}

\begin{document}

\hypertarget{ux95eeux9898ux4e00ux6a21ux578bux5206ux6790ux62a5ux544a}{%
\section{问题一模型分析报告}\label{ux95eeux9898ux4e00ux6a21ux578bux5206ux6790ux62a5ux544a}}

\hypertarget{ux6570ux636eux4e0eux5efaux6a21ux8bbeux7f6e}{%
\subsection{1.
数据与建模设置}\label{ux6570ux636eux4e0eux5efaux6a21ux8bbeux7f6e}}

本次建模使用 2025 年全年数据作为训练集，目标变量为出水浊度
\texttt{NTU}。验证集使用2026年1月；

训练数据在建模前构造了 2、4、6 小时时滞变量，涉及变量为：

\begin{longtable}[]{@{}ll@{}}
\toprule
变量 & 时滞 \\
\midrule
\endhead
R/W NTU & 2h, 4h, 6h \\
FILT. NTU & 2h, 4h, 6h \\
ALUM & 2h, 4h, 6h \\
\bottomrule
\end{longtable}

最终样本量如下：

\begin{longtable}[]{@{}llr@{}}
\toprule
数据集 & 时间范围 & 样本量 \\
\midrule
\endhead
训练集 & 2025-01-01 07:00 到 2025-12-31 23:00 & 4377 \\
验证集 & 2026-01-01 07:00 到 2026-01-31 23:00 & 357 \\
\bottomrule
\end{longtable}

\hypertarget{ux6a21ux578bux7ed3ux679c}{%
\subsection{2. 模型结果}\label{ux6a21ux578bux7ed3ux679c}}

本次比较了随机森林、XGBoost 和两者简单平均集成模型。评价指标包括
MAE、RMSE 和 R2。

\begin{longtable}[]{@{}llrrr@{}}
\toprule
模型 & 数据集 & MAE & RMSE & R2 \\
\midrule
\endhead
随机森林 & 训练集 & 0.1313 & 0.1997 & 0.8534 \\
随机森林 & 验证集 & 0.1439 & 0.1677 & 0.3649 \\
XGBoost & 训练集 & 0.1078 & 0.1503 & 0.9170 \\
XGBoost & 验证集 & 0.1609 & 0.1909 & 0.1773 \\
平均集成 & 训练集 & 0.1180 & 0.1698 & 0.8940 \\
平均集成 & 验证集 & 0.1505 & 0.1764 & 0.2973 \\
\bottomrule
\end{longtable}

从验证集表现看，随机森林单模型的 MAE、RMSE 和 R2 均优于 XGBoost
与平均集成。XGBoost
在训练集上拟合更强，但验证集表现下降更明显，说明其泛化能力弱于当前调参后的随机森林。

平均集成模型相较 XGBoost
有明显改善，但没有超过随机森林单模型。因此，在当前特征和参数设置下，若以
2026 年 1 月验证集为准，随机森林是更优选择。

\hypertarget{ux5173ux952eux53d8ux91cfux5206ux6790}{%
\subsection{3.
关键变量分析}\label{ux5173ux952eux53d8ux91cfux5206ux6790}}

XGBoost 特征重要性排名靠前的变量如下：

\begin{longtable}[]{@{}rlr@{}}
\toprule
排名 & 特征 & 重要性 \\
\midrule
\endhead
1 & FILT. NTU & 0.3460 \\
2 & FILT. NTU\_LAG\_2H & 0.2991 \\
3 & R/W FLOW & 0.0720 \\
4 & R/W NTU\_LAG\_6H & 0.0445 \\
5 & FILT. NTU\_LAG\_4H & 0.0356 \\
6 & R/W NTU & 0.0326 \\
7 & RIVER LEVEL & 0.0302 \\
8 & R/W NTU\_LAG\_2H & 0.0220 \\
9 & FILT. NTU\_LAG\_6H & 0.0191 \\
10 & ALUM & 0.0170 \\
\bottomrule
\end{longtable}

结果表明，\texttt{FILT.\ NTU} 及其短时滞后项是预测出水 \texttt{NTU}
的核心变量。这符合工艺逻辑：过滤后浊度与最终出水浊度最直接相关。\texttt{R/W\ FLOW}、\texttt{R/W\ NTU}
和河流水位也提供了原水状态与运行负荷信息。

\hypertarget{ux5e74-2-ux6708ux6307ux5b9aux65e5ux671fux9884ux6d4b}{%
\subsection{4. 2026 年 2
月指定日期预测}\label{ux5e74-2-ux6708ux6307ux5b9aux65e5ux671fux9884ux6d4b}}

以下结果为随机森林和 XGBoost 的简单平均集成预测。

\begin{longtable}[]{@{}lrrr@{}}
\toprule
日期 & 日均预测 NTU & 最小预测 NTU & 最大预测 NTU \\
\midrule
\endhead
2026-02-01 & 0.3794 & 0.3660 & 0.3979 \\
2026-02-10 & 0.4511 & 0.4176 & 0.4856 \\
2026-02-20 & 0.3634 & 0.3366 & 0.3802 \\
\bottomrule
\end{longtable}

逐时预测结果已保存至：

\texttt{Question1/outputs/xgboost\_ensemble/feb\_1\_10\_20\_predictions.csv}

从预测结果看，2 月 10 日整体预测值最高，日均预测 \texttt{NTU} 为
0.4511；2 月 20 日最低，日均预测 \texttt{NTU} 为
0.3634。三个日期的预测值均处于较低浊度区间，没有出现明显异常高值。

\hypertarget{ux7ed3ux8bba}{%
\subsection{5. 结论}\label{ux7ed3ux8bba}}

\begin{enumerate}
\def\labelenumi{\arabic{enumi}.}
\tightlist
\item
  在当前验证方式下，随机森林模型验证集表现最好，验证集 MAE 为
  0.1439，RMSE 为 0.1677，R2 为 0.3649。
\item
  XGBoost
  在训练集上拟合能力更强，但验证集效果较弱，说明存在一定过拟合或跨期泛化不足。
\item
  平均集成模型提高了 XGBoost 的稳定性，但未超过随机森林单模型。
\item
  影响出水浊度的关键变量主要是
  \texttt{FILT.\ NTU}、\texttt{FILT.\ NTU\_LAG\_2H}、\texttt{R/W\ FLOW}、\texttt{R/W\ NTU}
  相关变量。
\item
  对 2026 年 2 月 1 日、10 日、20 日的预测显示，2 月 10
  日出水浊度水平相对更高，2 月 20 日相对更低。
\end{enumerate}

\end{document}
